\section{Related Work and Theoretical Foundations}

\subsection{Solubility Prediction Methods}

Solubility prediction across media and temperatures has evolved over several decades into three broad families: empirical structure--property relationships, physicochemical thermodynamic models, and modern data-driven machine-learning models.

In the empirical QSPR family, the canonical examples are Delaney's ESOL model\cite{delaney2004} and Yalkowsky's General Solubility Equation (GSE). ESOL estimates aqueous solubility from only a handful of two-dimensional topological descriptors---the octanol--water partition coefficient ($\mathrm{MolLogP}$), molecular weight ($\mathrm{MW}$), the number of rotatable bonds, and the aromatic-ring proportion. These empirical relationships, however, are confined to aqueous media at room temperature (298.15 K) and cannot be carried over, mechanistically, to multi-solvent systems or continuous temperature changes.

Physicochemical and activity-coefficient thermodynamics, by contrast, centre on phase equilibria, giving rise to the UNIFAC group-contribution method, COSMO-RS (based on the screening-charge distribution of a conductor-like reference for real solvents), and the general empirical activity model NRTL-SAC. These physical models incorporate temperature and multi-component phase equilibria naturally, but their practical adoption is held back by a heavy dependence on the solid-state melting thermodynamics of the pure solute---the experimental melting point $T_m$ and the enthalpy of fusion $\Delta H_{\mathrm{fus}}$. Most novel candidate scaffolds have never been synthesised and crystallised in early discovery, so thermodynamic models cannot be applied directly in high-throughput screening of unseen parents.

More recently, hybrid thermodynamics and modern machine learning have tried to bridge first-principles reasoning and statistical fitting. Vermeire et al.\ combine rigorous thermodynamic derivation with machine-learned solvation free energies and enthalpies to extrapolate solubility over wide temperature and solvent ranges\cite{vermeire2022}; Attia et al.'s FastSolv contrasts message-passing neural networks against fixed descriptors for cross-molecule extrapolation on a large multi-solvent dataset\cite{fastsolv2025}. When only SMILES strings and operating temperatures are available, building an efficient regressor that is transparent in its physical prior yet generalises in a purely data-driven way remains the frontier at the intersection of computational chemistry and machine learning.

\subsection{Molecular Descriptors and Interaction Features}

For tabular learners, the design of the input features sets the ceiling on what the model can represent.

At the level of molecular representation, cheminformatics toolkits such as RDKit extract one- and two-dimensional physicochemical descriptors in a systematic way, covering topological indices, electronic distributions, surface areas, and partial-charge assignments; these features carry clear physical meaning and suit tabular learners naturally. Structural fingerprints (Morgan/ECFP4) capture local substructure fragments\cite{ye2021}, while graph neural networks learn continuous embeddings directly on the atom--bond graph\cite{fastsolv2025}. In extrapolation tasks split by parent, high-dimensional sparse fingerprints and deep graph networks tend to overfit the specific scaffolds of the training set, whereas dense physicochemical descriptors usually generalise better across parents.

Multi-solvent dissolution thermodynamics further demands representations that go beyond single-molecule properties. Dissolution reflects the competition among solute--solute, solvent--solvent, and solute--solvent interactions, microscopically following the like-dissolves-like rule. If the drug and solvent vectors are merely concatenated into $[D; S; T]$, an axis-aligned tree model needs many split depths to reach the cross-molecular nonlinear combinations. Boobier et al.\ note that explicitly forming the absolute differences of matching properties, $|D-S|$ (polarity mismatch $|\Delta\mathrm{TPSA}|$, lipophilicity mismatch), and the Hadamard products $D\odot S$ (co-polarity and hydrogen-bond pairing strength) injects physicochemical priors directly into the learner and shrinks the feature-search space considerably\cite{boobier2020}.

\subsection{Structure--Temperature Modelling}

In physical thermodynamics, the classical van 't Hoff equation relates the solid--liquid equilibrium log-solubility to the inverse absolute temperature in a linear fashion, with the slope controlled by the apparent enthalpy of dissolution $\Delta H_{\mathrm{sol}}$. Over the moderate operating window relevant in practice, a first-order local Taylor expansion implies that solubility responds to temperature in a quasi-linear or monotonic way.

Traditional machine-learning methods tend to feed temperature as a single column alongside hundreds of structural attributes, where its signal is easily diluted and ignored; or, because the sampling conditions are non-random, the learner picks up spurious temperature shortcuts. The SOLAR-H architecture proposed here borrows the fixed-effects idea from econometrics and cancels the structural constant inside each drug--solvent combination, so that the true local temperature slope is recovered from pure within-group variation---uniting physical mechanism and statistical learning.

\subsection{Ensembles and Tabular Foundation Models}

On tabular regression, tree ensembles---random forests, ExtraTrees, CatBoost\cite{catboost2018}---have long held the mainstream. Extremely Randomized Trees (ExtraTrees\cite{extratrees2006}), in particular, randomise both the candidate features and the split thresholds at each node; relative to a standard random forest they smooth the regression staircase more effectively, lower the prediction variance, and resist collinearity noise in high-dimensional descriptor space.

A parallel breakthrough comes from tabular deep learning. The TabPFN v2 model of Hollmann et al., published in \textit{Nature}, demonstrates the strength of tabular foundation models\cite{tabpfn2025}. Pretrained offline on millions of synthetic causal-prior datasets, TabPFN takes the entire labelled training set as context at inference time (in-context learning) and implicitly performs efficient Bayesian posterior-mean inference. Its smooth metric-space inductive bias is highly complementary, in learning mechanism, to the axis-aligned hyperplane splits of ExtraTrees---an argument that motivates blending the two.

\subsection{Evaluation, Interpretability, and Causal Inference}

On methodology and mechanism, recent work has codified several important norms. On cross-sample leakage, Kapoor and Narayanan document the widespread overestimation in applied machine-learning science\cite{leakage2023}, and Llompart et al.\ stress how dataset overlap and chemical-structure similarity mislead the evaluation of solubility models\cite{solubility2024}. A fully isolated holdout by parent scaffold (GroupKFold by parent) has thus become the standard for measuring true generalisation. Nested cross-validation goes one step further, sealing the hyperparameter search and weight selection entirely inside the outer training folds to avoid second-order leakage\cite{varma2006,cawley2010}. On decisions and causality, the TreeSHAP algorithm of Lundberg et al.\ computes exact Shapley-based local and global feature contributions for tree ensembles in polynomial time\cite{treeshap2020}; extracting physical mechanism from uncontrolled observational data, in turn, requires causal-inference theory\cite{pearl2009}, which here operates by de-meaning within combinations to block the main confounding paths, combined with the Cinelli--Hazlett sensitivity framework to bound the robustness against unmeasured environmental confounders\cite{sensitivity2020}.

\subsection{Positioning}

Prior work has explored a rich set of molecular representations and algorithmic designs for solubility prediction, yet in multi-solvent, variable-temperature, and unseen-parent settings, deep problems remain: sampling imbalance is left uncorrected, evaluation leakage is uncontrolled, and the temperature effect is confounded.

This work targets these pain points directly. On a strictly leakage-free parent-nested evaluation, it introduces SOLAR-H, a decomposition that debiases through physical insight, and ExtraTrees--TabPFN, a heterogeneous blend of 869-dimensional explicit interactions, parent weighting, and the TabPFN foundation model. Through extensive ablation, TreeSHAP attribution, fixed-effects causal identification, and an interactive deployed system, the work offers an end-to-end solution for predicting molecular properties under complex conditions.
